% ============================================================
\section{Related Work}
\label{sec:related}
% ============================================================

\textbf{Traditional IoT IDS.}
Statistical methods (Z-score, IQR) and Isolation
Forest~\cite{liu2012isolation} flag individual flow features deviating
from learned normals; ensemble approaches~\cite{garcia2014empirical}
improve precision.  All share a weakness: calibration on known,
high-magnitude attacks and failure on stealthy perturbations.
We use the \dataset{} dataset~\cite{neto2023ciciot} throughout.

\textbf{Deep-learning anomaly detection.}
USAD~\cite{audibert2020usad} uses adversarially trained dual-autoencoders.
TranAD~\cite{tuli2022tranad} applies context-aware Transformers.
The Anomaly Transformer~\cite{xu2022anomaly} introduces association
discrepancy for anomalous timestep detection.
PatchTST~\cite{nie2023patchtst} captures long-range dependencies via
patch-based Transformers.
We include all four as baselines.

\textbf{Generative augmentation for security.}
GANs~\cite{lin2018idsgan} suffer mode collapse; TimeGAN~\cite{yoon2019timegan}
improves temporal fidelity; Diffusion-TS~\cite{zhang2024diffusionts}
achieves state-of-the-art quality for time-series generation.
We use Diffusion-TS as a backbone for generating \emph{learned}
negatives: we train a Diffusion-TS model on benign IoT traffic, then
apply the same analytical perturbation pipeline to generate
protocol-constrained stealth-controlled negatives.
This provides a controlled comparison between analytical (closed-form)
and learned (generative-model) negatives, revealing that negative
quality mediates catastrophic forgetting at scale.

\textbf{Adversarial training.}
FGSM~\cite{goodfellow2015adversarial} and PGD~\cite{madry2018pgd} generate
adversarial perturbations at test time.  We generate adversarial
\emph{training data} via closed-form perturbations targeting the
traffic-generation threat model, with protocol constraints ensuring
domain validity.

\textbf{Time-series foundation models.}
Moirai~\cite{woo2024moirai}, trained on 27B+ time points, supports
probabilistic multivariate forecasting.
Its confidence-interval output maps naturally to anomaly scoring.
MOMENT~\cite{goswami2024moment} and TimesFM~\cite{das2024timeseries}
offer complementary architectures.

\textbf{Catastrophic forgetting and continual learning.}
Catastrophic forgetting---the phenomenon where fine-tuning overwrites
previously learned representations---is well-studied in continual
learning.  EWC~\cite{kirkpatrick2017ewc} penalises changes to important
weights; PackNet~\cite{mallya2018packnet} prunes and freezes network
subsets; Progressive Neural Networks~\cite{rusu2016progressive} add
lateral connections to frozen columns.  In the foundation-model era,
LoRA~\cite{hu2022lora} and adapters~\cite{houlsby2019adapters} avoid
forgetting by training low-rank or bottleneck modules while freezing
the base model.  We do not compare against EWC or PackNet empirically; our encoder-freezing
and LoRA experiments (Section~\ref{sec:experiments:scaling}) are diagnostic
tools to confirm the forgetting mechanism, not proposed mitigations.

\textbf{Contrastive learning for anomaly detection.}
CSI~\cite{tack2020csi} uses contrastive learning with distributional
shift detection for out-of-distribution samples.
NeuTraL-AD~\cite{qiu2021neutral} learns neural transformations for
anomaly scoring via contrastive objectives.
SSD~\cite{sehwag2021ssd} combines self-supervised representations with
Mahalanobis-distance outlier detection.
These methods operate in the image domain; our work extends the
contrastive anomaly detection principle to multivariate time series
with domain-specific (IoT protocol) constraints.
Supervised contrastive learning~\cite{khosla2020supcon} has been applied
to tabular~\cite{sohn2020fixmatch} and time-series~\cite{eldele2021ts2vec}
anomaly detection.  Our CombinedLoss adapts~\cite{reiss2021panda}'s
reconstruction$+$contrastive principle to IoT time series.
Our stealth-controlled negatives are analytically specified near the benign boundary,
preventing the SupCon collapse observed with high-magnitude attacks.
